\section{Emerging Practices}
\label{sec:emerging}

The practices above adapt established methods. This section describes nine
that are newer: some were published in the last two years, some have been
tried first in HRI, and some are proposals that follow from the framework and
have not, to our knowledge, been tested. We state for each what it can
establish and where its limits lie.

\subsection{Ceiling-Referenced Agreement and Tests of Substitution}
\label{sec:equivalence}

Reporting model--human agreement next to human--human agreement
(\bpr{bp:human}, \bpr{bp:coder}) invites a summary: the model reaches some
fraction of the human ceiling. The summary is useful descriptively, but a
ratio of two noisy coefficients is unstable in small samples and does not
answer the question usually being asked, which is whether the model may be
used in place of a human coder or rater.

The \emph{alternative annotator test} addresses that question directly
\cite{calderon2025alternative}. With a small group of human annotators and a
modest number of items annotated by all of them, it asks, for each annotator
in turn, whether the model agrees with the remaining annotators at least as
well as the excluded annotator does, allowing a stated margin for the
model's lower cost; the model is accepted as an alternative if it wins this
comparison for at least half of the annotators. The procedure fits HRI's
circumstances: its authors recommend at least three annotators and 50 to 100
items, which is the multiply annotated subset that \bpr{bp:human} already
calls for, and it yields a justified decision in place of an appeal to a
conventional threshold. Variants extend the idea to the common case in which
different items were rated by different people, by testing whether the
agreement of panels that include the model is non-inferior to that of
all-human panels \cite{omar2026altratest}, and to equivalence testing on
Krippendorff's $\alpha$ with the model swapped into the group of human
annotators \cite{he2025can}. A related proposal expresses a judge's
agreement with people as a standardised distance from human--human
agreement, which distinguishes judges that agree with people about as much
as people agree with each other from those that are implausibly consistent
\cite{han2025judges}.

Two cautions apply to any ceiling-referenced figure. A single aggregate can
flatter the model: when many items have uncertain human labels,
model--human agreement can appear equal to or better than human--human
agreement, while on the items where people are consistent it falls well
below, so agreement should be stratified by how much the human raters
agreed \cite{elangovan2025beyond}. And a model should be scored against one
rater at a time and the results averaged, not against the panel's majority
vote; the comparison can then be expressed as the number of human raters
whose pooled judgement the model matches \cite{resnick2021principled}.

More generally, any substitution claim is a claim of equivalence or
non-inferiority and should be tested as one. The two one-sided tests
procedure asks the researcher to state in advance the smallest difference
that would matter and then tests whether the observed difference lies within
those bounds \cite{lakens2017equivalence,lakens2018equivalence}. The logic
applies to a model compared with a human coder, to model-authored responses
compared with human-authored ones (\bpr{bp:humanref}), to an autonomous robot
compared with a wizard-controlled one, and to claims that an outcome does not
differ between conditions. It has the salutary effect of making authors say
how much worse than a human the model is allowed to be.

\subsection{Generalisability Designs: Runs and Prompts as Raters}
\label{sec:gtheory}

Classical reliability treats one source of error at a time. An evaluation of
a generative model has several at once: the items sampled, the human raters,
the model's stochastic runs, the wording of the prompt, and the model
version. Generalisability theory was developed for exactly this situation in
behavioural measurement
\cite{shavelson1991generalizability,brennan2001generalizability}, and
recent guidance extends it to incomplete designs in which raters only
partly overlap across items \cite{tenhove2024updated}. In a
generalisability study, items are crossed with the other facets (each of a
sample of items is scored under several prompt variants, several runs each,
and by several human raters), and the variance of the scores is decomposed
into components attributable to each facet and their interactions. A
decision study then uses those components to compute how reliable a score
would be under any proposed design.

Applied to a model in a Judge or Annotator role, this answers practical
questions that are otherwise settled by habit. If run-to-run variance is
small and prompt-to-prompt variance is large, averaging over five runs of one
prompt buys nothing and averaging over three prompts buys a great deal. If
the item-by-rater interaction dominates, the construct is one on which people
differ and no change to the model will raise agreement. The design also
gives a principled meaning to treating the model as a rater: its runs and
prompts are facets on the same footing as human raters, and their
contribution to error is estimated, not assumed. The cost is a fully crossed
subset of perhaps fifty to a hundred items, which is within reach of most
studies. We are not aware of an HRI study that has reported such a
decomposition, and we suggest it as the natural quantitative form of
\bpr{bp:consistency}.

\subsection{Disagreement-Aware Evaluation}
\label{sec:perspectivist}

Where human raters legitimately disagree (\bpr{bp:human}), evaluating a model
against a majority label measures its agreement with a fiction. The
perspectivist programme in NLP proposes keeping the individual labels and
evaluating against their distribution
\cite{plank2022problem,cabitza2023toward,uma2021learning}. For HRI this suits
constructs such as social appropriateness, politeness, and comfort, which
vary with culture, age, and familiarity with robots. The datasets needed
already exist in part: collections of household scenes in which thousands of
annotators scored the appropriateness of robot actions and explained their
scores retain the spread of judgements for each scene, and GRACE uses that
spread directly, separating scenes on which annotators agree from those on
which they do not and treating the two differently \cite{dogan2025grace}.

Evaluation against a majority label is also technically unsound in
this situation: calibration measured against the human majority is
ill-defined when people inherently disagree \cite{baan2022stop}.

Three designs follow. A Perceiver or Judge can be scored by how closely the
distribution of its labels over repeated runs or personas matches the
distribution of human labels for the same item, for instance by a divergence
measure, with the caution that variability across a model's samples is not
the same thing as variability across a population and tends to be narrower
\cite{bisbee2024synthetic,wang2025large}. Agreement can be reported within
rater subgroups, so that it is visible when a model tracks the judgements of
one group and not another. And a Speaker's or Actor's behaviour can be
evaluated by the proportion of raters who find it acceptable, with the lower
tail of that distribution, the people for whom it was not acceptable,
reported explicitly, since in deployment those people are users too.

\subsection{Frozen-Interaction Replay}
\label{sec:replay}

A hosted model that has been withdrawn cannot be re-run, and a study
conducted with it cannot be replicated in the usual sense. What can be
preserved is the set of situations the model faced. We propose that studies
archive, for every model call made during an interaction, the complete
context supplied to the model (prompt, dialogue history, retrieved memory,
sensor descriptions or images) and the output returned. This archive supports
three uses.

\emph{Drift quantification.} When a model is updated, the archived contexts
are replayed through the new version and the outputs compared with the
originals on the study's rubric, by the study's raters or its validated
judge. The result is a direct estimate of how much the robot's behaviour
would have changed, in the units the study cares about, which is the
information a reader needs to decide whether the original finding still
applies. Given the documented scale of behavioural change between versions of
the same commercial model \cite{chen2024how}, this is not a hypothetical
concern.

\emph{Regression testing.} Contexts in which the original model failed, or
in which a supervisor intervened, become a test suite. The idea of testing
specific behavioural capabilities with targeted cases, instead of relying on
aggregate held-out accuracy, is established in NLP \cite{ribeiro2020beyond};
an HRI laboratory's own failures are the most relevant cases it will ever
have.

\emph{Cross-laboratory comparison.} With consent and de-identification,
archived contexts can be shared, so that different models and prompts are
compared on identical situations drawn from real interactions with real
populations, which no public benchmark currently offers.

Replay of recorded field interactions has been used in HRI to evaluate a
perception component: footage from a multi-day Wizard-of-Oz deployment in a
café was replayed to a vision-language model to test whether it could read
customers' readiness to be served \cite{bu2025using}. The same logic
applies to any component whose input can be recorded.

The limit of replay is that it is open-loop. It shows what a new model would
have said at a given point in an old conversation, not how the conversation
would then have gone. Replay therefore supports component-level and not
interaction-level claims. For action models the analogue is evaluation in
simulation built to correlate with real outcomes
\cite{li2025evaluating,jain2026polaris}, or in learned world models
validated against real evaluations \cite{geminirobotics2025evaluating},
with a small frozen set of physical scenarios as an anchor, and the same
caveat applies with greater force because physical interaction is
closed-loop at every time step.

\subsection{Supervised Autonomy as an Evaluation Instrument}
\label{sec:lloz}

Between the Wizard-of-Oz study, in which a person produces the robot's
behaviour, and the autonomous study, in which a model does, lies an
arrangement that deserves to be used deliberately: the model proposes and a
trained supervisor approves, edits, or replaces each proposal before it is
delivered. The arrangement is usually adopted for safety, and it is the
ethically appropriate default for early Tier~III studies. Its value as an
evaluation instrument is less often noticed. Every decision by the supervisor
is an expert judgement of a model output, made in context, at the moment of
use. The proportion of proposals approved unchanged, the size and kind of
edits, and the reasons for replacement constitute a component-level
evaluation that is collected as a by-product of running the study, on
exactly the distribution of situations that the deployment will meet.

Reporting should follow the spirit of the Wizard-of-Oz guidelines
\cite{riek2012wizard}: who the supervisor was and how they were trained, what
they could see, the time pressure under which they decided (since an
approval given within a conversational pause is a weaker endorsement than
one given at leisure), and whether participants knew. A subset of
proposals should be independently rated afterwards so that the supervisor's
own reliability is known. The sequence from wizard-controlled to supervised
to autonomous operation, with the intervention rate reported at each stage,
gives a study programme a graded path to autonomy in which each step is
licensed by evidence from the last.

\subsection{Embodied Red-Teaming and Psychosocial Safety}
\label{sec:redteam}

Adversarial testing before release is standard for language models and is
recommended by risk-management frameworks for generative AI
\cite{nist2024artificial}. Robots change what is at stake. Jailbreaking
attacks devised for chatbots have been adapted to LLM-controlled robots and
have elicited harmful physical actions from research and commercial
platforms, in some settings with attack success rates of 100\%
\cite{robey2025jailbreaking}, including through ordinary voice interaction
\cite{zhang2025badrobot}; in an interactive simulator the most
safety-conscious of the agents tested rejected only one in ten detailed
hazardous tasks \cite{yin2024safeagentbench}; LLMs evaluated in
HRI task scenarios have accepted instructions that were dangerous, unlawful,
or discriminatory \cite{hundt2025llm}; and benchmarks of semantic safety for
embodied models have been constructed together with automatically generated
robot ``constitutions'', the best of which still left roughly one case in
six misaligned \cite{sermanet2025generating}. For robots in social
and supportive roles, the hazards are psychosocial as much as physical:
manipulation, the fostering of dependence, harmful advice, and breaches of
confidence, for which multi-dimensional safety evaluation has begun to be
developed \cite{luo2026dialogguard}. Defences can be evaluated in the same
way: a guardrail architecture reduced the execution of unsafe plans under
worst-case jailbreak attacks from over 92\% to below 3\%
\cite{ravichandran2026safety}, and stress tests need not put hardware or
people at risk, since hard instructions can be generated automatically
\cite{karnik2024embodied} and degradation under off-nominal conditions
predicted from edited observations before any physical trial
\cite{majumdar2025predictive}.

A first study of this kind in HRI invited 21 students to try to make an
LLM-driven social robot violate principles concerning attachment, freedom,
and empathy, and identified five manipulation techniques
\cite{abbo2025can}.

We recommend that Tier~III studies include a documented red-team exercise
before participants are exposed, and that Tier~II studies with unreviewed
user-facing output include a lighter one. The report should state who
performed it, which categories of hazard were probed, how many attempts were
made, how many succeeded, and what was changed as a result. The categories
should include those specific to embodiment: instructions that are harmless
as text and harmful as action; instructions from a bystander who is not the
user; attempts to have the robot disclose what another user told it; and
content in the robot's field of view or hearing that the model may treat as
an instruction. People with relevant expertise, including clinicians and
safeguarding professionals where the population warrants, should take part,
as should members of the population itself where that is feasible.

\subsection{Persona and Identity Stability}
\label{sec:persona}

A robot's character is part of its design, and with a language model it is
specified in a prompt and maintained, or not, by the model across a
conversation. HRI experts consulted on prompt design identified the
legibility of the robot's personality, users' ability to adapt it, and
safety, deception, and governance as their main concerns \cite{ashok2026why}.
Drift is measurable and can be fast: adherence to a system-prompt persona
declined significantly within eight rounds of conversation in two widely
used models \cite{li2024measuring}. Stability can be measured. Probes of the persona's stated traits, knowledge,
and boundaries are inserted at intervals in long simulated or replayed
conversations, and the consistency of the answers is scored; the same
conversations show whether the robot's account of its own nature and
capabilities stays truthful. Work on multi-turn simulation shows that persona drift and echoing between
agents arise in long generated dialogues and that how the conversation
history is presented to each agent substantially affects them
\cite{luo2026spasm}, which implies
that persona stability is a property of the whole configuration and has to
be tested at the conversation lengths that the deployment will reach. We
caution against measuring a robot's ``personality'' by administering human
personality inventories to the model: the psychometric properties of such
inventories do not automatically transfer to language models
\cite{suhr2025challenging}, and what matters for HRI is the persona that
people perceive, which is an interaction-level question. Studies of
LLM-driven robot personalities that verify the perceived persona with
participants before using it as an experimental manipulation show the
appropriate order of work \cite{becker2025robots,jafari2026evaluating}.

\subsection{Reference-Corpus Alignment}
\label{sec:alignment}

A robot that offers support, tutoring, or advice is implicitly compared with
the human professionals who do the same. Matched human responses
(\bpr{bp:humanref}) are the direct comparison, but they are expensive.
Reference-corpus alignment is a scalable complement: the robot's
conversations are compared with an existing corpus of professional practice
in a shared embedding space. In the first application to HRI, disclosures
made to an LLM-driven robot during a wellbeing intervention and the robot's
responses were compared with a corpus of counselling exchanges between
help-seekers and licensed therapists. Themes were derived by clustering each
corpus; each item was then assigned to the nearest theme of the other corpus
and counted as fitting if it lay within a distance threshold derived from
that theme's own dispersion. About 91\% of disclosures to the robot fitted
therapy themes, about 28\% of therapy disclosures fitted robot-conversation
themes, and both disclosures and responses were compared for semantic
similarity within matched themes \cite{chiang2025talk}.

The method establishes something that no rubric-based audit of a sample can:
where, in the space of what professionals deal with, the robot's
conversations fall, and which regions of professional practice they never
enter. The asymmetry in the figures above is such a finding. Its limits are
equally clear. Proximity in embedding space is not appropriateness; a
response can be topically aligned with therapists' responses and still be
wrong for the person. Results depend on the embedding model and require the
checks of \bpr{bp:embed}. Dispersion-based thresholds give diffuse reference
themes a larger capture radius, which the original authors note. And the
reference corpus must be examined for comparability: in the case described,
written exchanges on a counselling platform were compared with spoken
conversations between students and a robot. Used with these cautions,
alongside and not instead of rated audits, reference-corpus alignment gives
HRI a way to situate generative robots relative to human practice at a scale
that hand coding cannot reach.

\subsection{Shared Infrastructure}
\label{sec:infrastructure}

Several of the practices in this article are costly for a single laboratory
and cheap for a community. Three kinds of shared resource would lower the
cost. The first is tooling that makes good reporting automatic: experiment
platforms that version prompts and configurations and log every model call
with its context, as open-source tools for LLM-based interaction experiments
already do for text agents \cite{laban2024lexi}. The second is shared,
multiply rated collections of HRI material (robot utterances in context,
clips of generated behaviour, videos of handovers) with which a laboratory
can validate a judge for a common construct without rating everything anew;
these would be to HRI what meta-evaluation datasets are to NLP. The third is
distributed evaluation, in which the same configuration is evaluated at
several sites, as robot learning has begun to do for generalist policies
\cite{atreya2025roboarena} and as psychology does in multi-site
replications; for interaction-level claims about generative robots, whose
behaviour differs for every participant, multi-site evidence is likely to be
the only route to conclusions that hold beyond one laboratory's population.
